%%%%%%%%%%%%%%%%%%%%%%%%%%%%%%%%%%%%%%%%%
% Twenty Seconds Resume/CV
% LaTeX Template
% Version 1.0 (30/6/20)
%
% Original author:
% Carmine Spagnuolo (cspagnuolo@unisa.it) and with major modifications by 
% Vel (vel@LaTeXTemplates.com), and Harsh (harsh.gadgil@gmail.com) 
% and Rougui (jamal.rougui@gmail.com) :
% + translate to french language
% + add profile section
% + fix some compilation problem and style
%
% License:
% The MIT License (see included LICENSE file)
%
%%%%%%%%%%%%%%%%%%%%%%%%%%%%%%%%%%%%%%%%%

%----------------------------------------------------------------------------------------
%	PACKAGES AND OTHER DOCUMENT CONFIGURATIONS
%----------------------------------------------------------------------------------------
\documentclass[letterpaper]{twentysecondcv} % a4paper for A4
%\newfontfamily{\FAA}{[FontAwesome.otf]}
% Command for printing skill overview bubbles
\usepackage{fontawesome}
\newcommand\skills{ 
	~
	\smartdiagram[bubble diagram]{
		\textbf{Expertise}\\\textbf{Logiciels},
		\textbf{~~Bonnes ~~}\\\textbf{~~pratiques~~},
		\textbf{~~Télécom~~},
		\textbf{~~~~~~Sécurité~~~~~~},
		\textbf{~~~Machine~~~}\\\textbf{~~~Learning~~~},
		\textbf{~~~Design pattern~~~},
		\textbf{~~~Optimisation~~~},
		\textbf{~~~Cloud~~~}
	}
}

\aboutme{\\
	{\large \textbf{15}} ans d'expérience professionnelle dans le métier d'ingénierie informatique, à l'écoute, bon communiquant et force de proposition.
	
	{\large \textbf{5}} ans d'expérience en R\&D me procurant une forte capacité de trouver des solution).
	%Passionné par l'innovation et la veille technique, toujours attiré par des projet innovant et "from skratch" faisant appel à des outils et Frameworks maintenable et libres.
}

% Programming skill bars
\programming{{Cloud $\textbullet$ AKS  $\textbullet$ SRE / 3}, {JEE $\textbullet$ SOA $\textbullet$ \large Angular / 3.5}, {Java $\textbullet$ Web $\textbullet$ Design pattern / 5}}

% Projects text
\education{
	\textbf{Doctorat, Sciences d'ingénieries}\\
	Université Nantes \& UM5R \\
	2004 - 2008 | Nantes, France
	
	\textbf{M2., Info\&Télécom} \\
	Université Mohammed V de Rabat (UM5R)\\
	2002 - 2004 | Rabat, Maroc
}

%----------------------------------------------------------------------------------------
%	 PERSONAL INFORMATION
%----------------------------------------------------------------------------------------
% If you don't need one or more of the below, just remove the content leaving the command, e.g. \cvnumberphone{}

\cvname{Jamal ROUGUI} % Your name
\cvjobtitle{Lead Tech / Concepteur} % Job
% title/career

\cvlinkedin{/in/jamal-rougui-45462624/}
\cvgithub{jerougui/jerougui.github.io}
\cvnumberphone{(+33) 6 60 444 936} % Phone number
\cvsite{credly.com/users/jamal-rougui/} % Personal website
\cvmail{jamal.rougui@gmail.com} % Email address

%----------------------------------------------------------------------------------------

\begin{document}
	
	\makeprofile % Print the sidebar
	
	%----------------------------------------------------------------------------------------
	%	 EXPERIENCE
	%----------------------------------------------------------------------------------------
	
	\section{Expériences}
	
	\begin{twenty} % Environment for a list with descriptions
		
		\twentyitem
		{à ce jour -}	
		{Mai 2023}
		{LeadTech / Concepteur}
		{\href{https://www.transilien.com/fr}{SNCF Voyageurs DSI Translien, Nantes}}
		{}
		{
			\begin{itemize}
				\item Définir une roadmap technologique pour traiter l'obsolescence, le décommissionnement et les migrations techniques
				\item capitalisation et rédaction des spécification technique pour répondre au EB
				\item Superviser l'implémentation de solutions techniques par les sous-traitants
				\item Support technique pour tous les acteurs de l'entité
				\item Veiller à l'application des meilleures pratiques et former les équipes 
				\item Veille technique et formation continue, participation à des salon tech                                                            
				
			\end{itemize}
		}
		
		\twentyitem
		{Mai 2023 -}
		{Mai 2022}
		{Lead Tech \& Aspiring architect}
		{\href{https://www.transilien.com/fr}{DSI SNCF TN, Capgemini Nantes}}
		{}
		{\begin{itemize}	
				\item Expertise technique en architecture orientée service dans un environnement cloud
				\item Conception et macro-chiffrage d'évolution d'un parc applicatif cloud en mode DevOps
				\item Veille technologique, audit du code, support MCO, astreint
				\item Point de contact BUILD et RUN avec les architectes afin d'améliorer les outils et vérifier la conformité des API
		\end{itemize}}
		\\
		\twentyitem
		{Avril 2022 -}
		{Jun 2015}
		{Lead Tech \& Agile master}
		{\href{https://www.thalesgroup.com/fr}{Thales SIX GTS France}}
		{}
		{\begin{itemize} 		
				\item Développement client lourd, autour des outils de tests et de validation en mode Feature Team.
				\item Sprint planning, revu, estimation, découpage, faisabilité
				\item développement et test, qualité, intégration continue
				\item force de proposition et amélioration continue
				\item Démo, review, rétrospective, report
		\end{itemize}}
		\\
		\twentyitem    	
		{Jun 2015 -}
		{Oct 2014}
		{Lead Tech \& Proxy Chef de Projet}
		{\href{https://www.insee.fr/fr/accueil}{INSEE, SII-Nantes}}
		{}
		{
			{\begin{itemize}
					\item Études de faisabilité et spécification formelle
					\item Migration IE to Firefox, et mise en place du SSO (JEE)
					\item Maintenance Adaptative et évolutive d'application (ASP)
					\item Maintenance Adaptative du Socle Java et Atelier de développement
			\end{itemize}}
		}
		\\
		\twentyitem   		
		{Jun 2015 -}
		{Avr 2013}
		{Développeur confirmé Java/Swing}
		{\href{}{HMX-HYPERVISION}}
		{}
		{
			\begin{itemize}
				\item Maintenance évolutive et corrective des logiciels vote électronique
				\item Développement des solutions
				\item Développement tests et packaging multi OS et support
			\end{itemize}
		}
		\\   
		\twentyitem
		{Sep 2014 -}
		{Mai 2014}
		{Développeur JEE}
		{\href{https://www.nexter-group.fr/}{NEXTER}}
		{}
		{
			\begin{itemize}
				\item Migration d'un prototype client lourd (java Swing) en client léger
				\item Conception d'architecture logicielle dans un environnement JEE
				\item Amélioration continue et métrique de la qualité du code
				\item Développement, test, optimisation des performances BDD.
			\end{itemize}
		}       
		\\
		\twentyitem
		{Nov 2012 -}
		{Oct 2011}
		{Développeur Java/C++}
		{\href{https://www.jobs.bouyguestelecom.fr/}{CDN - Bouygues Telecom}}
		{}
		{
			\begin{itemize}
				\item Maintenance adaptative et évolutive intervention dans toutes les étapes du cycle en V (spécification, développement, tests, support et maintenance) du SI Facturation et comptabilité auxiliaire
				\item Rédaction de spécifications techniques
				\item Développement Java/C++/PLSQL/KSH
				\item Packaging et livraison	
			\end{itemize}
		}
		
		\\
		\twentyitem
		{Oct 2011 -}
		{Mai 2011}
		{Développeur Java}
		{\href{}{E-Map, École Polytech'Nantes}}
		{}
		{
			\begin{itemize}
				\item Développement d'un moteur de recherche sémantique des données
				%\item Implémentation des solutions de traitement des grands volumes de données, classification et recommandation	
			\end{itemize}
		}
		
		\\
		%	\twentyitem
		%	{Nov 2009 -}
		%	{Aoû 2010}
		%	{Attaché Temporaire Enseignement et Recherche (ATER)}
		%	{\href{http://www.lina.univ-nantes.fr/}{LINA}}
		%	{}
		%	{
			%	\begin{itemize}
				%	\item Enseignement de l'informatique au niveau L1/L2
				%	\item Algorithmique et programmation, Bases de données
				%	\end{itemize}
			%	}
		%\twentyitem{<dates>}{<title>}{<location>}{<description>}
	\end{twenty}

%----------------------------------------------------------------------------------------
%	 RESEARCH
%----------------------------------------------------------------------------------------
\section{Recherche}
\begin{twenty}
	\twentyitem
    	{2004 - 2009}
		{}
        {R\&D. en Informatique}
        {\href{https://www.univ-nantes.fr/}{ESIGETEL \& Université de Nantes}}
        {}
	       {
	        \textbf{PostDoc}: Reconnaissance des sons de vie courante en environnement intelligent pour la santé - Robot CompanionAble. \\
	       	\textbf{Thèse}: Indexation de documents audio: cas des gros volumes de données.
	        {
	        %\begin{itemize}
	        %\item 
			%\end{itemize}
			}
	       }
  
\end{twenty}

\section{Publications}
\begin{thebibliography}{9}
	\bibitem{Rougui2007}
	J.E. Rougui, M. Gelgon, D. Aboutajdine, N. Mouaddib, M. Rziza.
	\emph{Organizing Gaussian mixture models into a tree for scaling up speaker retrieval}.
	Pattern Recognition Letters, In Press, Corrected Proof, Available online 19 March 2007.
	\url{http://cat.inist.fr/?aModele=afficheN&cpsidt=18887914}.
	
	\bibitem{Rougui2009a}
	J.E. Rougui, D. Istrate, W. Souidene, M. Opitz, M. Riemann.
	\emph{Audio based surveillance for cognitive assistance using a CMT microphone within socially assistive technology}.
	EMBC2009, September 2-6, Minneapolis, USA, 2009.
	Papier co-écrit avec AKG.
	
	\bibitem{Rougui2009b}
	J.E. Rougui, D. Istrate, W. Souidene.
	\emph{Audio Sound Event Identification for distress situations and context awareness}.
	EMBC2009, September 2-6, Minneapolis, USA, 2009.
	
	\bibitem{Rougui2006a}
	J.E. Rougui, M. Gelgon, M. Rziza, J. Martinez, D. Aboutajdine.
	\emph{Hierarchical organization of a set of gaussian mixture speaker models for scaling up indexing and retrieval in audio documents}.
	ACM Symposium on Applied Computing (ACM SAC'06), Multimedia and Visualization (MV) track, p. 1369-, 2006.
	
	\bibitem{Rougui2006b}
	J.E. Rougui, M. Rziza, D. Aboutajdine, M. Gelgon, J. Martinez.
	\emph{Fast incremental clustering of Gaussian mixture speaker models for scaling up retrieval in on-line broadcast}.
	IEEE International Conference on Acoustics, Speech, and Signal Processing (ICASSP '06), MMPSP-P3.6(V-521-524), Toulouse, France, 2006.
	
	\bibitem{Bzivin2003}
	J. Bzivin, G. Dup, F. Jouault, G. Pitette, J.E. Rougui.
	\emph{First experiments with the atl model transformation language: transforming xslt into xquery}. OOPSLA Workshop, 2003.

\end{thebibliography}
\end{document} 
